% Despliegue de los frontends
\subsection{Despliegue de los frontends}\label{sec:7_2_frontends}

\subsubsection{Ubicación del build web en dist/browser y asignación de permisos (755/644)}
Cada frontend tiene su carpeta en \nolinkurl{/apps/aplicaciones/defensoria} con su \texttt{Dockerfile}, su configuración de nginx, su script y, en \texttt{dist/browser}, los archivos compilados (sección \ref{sec:4_4_transferencia_vps}):
\begin{table}[H]
\centering
\caption{Carpetas de los frontends en el VPS Frontend}
\footnotesize
\begin{tabular}{p{4.4cm}p{3.6cm}p{6.4cm}}
\toprule
\textbf{Carpeta} & \textbf{Contenedor} & \textbf{Script} \\
\midrule
\texttt{front/} & defensoria-web & \nolinkurl{podman-compose-front.sh} \\
\texttt{front-admin/} & admin-web & \nolinkurl{podman-compose-front-admin.sh} \\
\texttt{front-revision/} & revision-web & \nolinkurl{podman-compose-front-revision.sh} \\
\texttt{front-primer-contacto/} & primer-contacto-web & \nolinkurl{podman-compose-front-primer-contacto.sh} \\
\texttt{front-subdefensoria/} & subdefensoria-web & \nolinkurl{podman-compose-front-subdefensoria.sh} \\
\bottomrule
\end{tabular}
\end{table}
Los archivos web son públicos (cualquiera los descarga desde el navegador), por eso llevan permiso 644 y sus carpetas 755; los scripts, 750:
\begin{lstlisting}
# F=/apps/aplicaciones/defensoria
# find $F -type d -exec chmod 755 {} +
# find $F -type f -exec chmod 644 {} +
# chmod 750 $F/*/podman-compose-front*.sh
# ls $F/*/dist/browser/index.html
\end{lstlisting}
El último comando debe listar los cinco \texttt{index.html}; si falta alguno, el frontend no se copió completo.
\captura{front-02-dist-browser}{Archivos compilados de los cinco frontends en dist/browser}

\subsubsection{Contenedores de los 5 frontends (nginx:alpine)}
Cada script construye la imagen con su Dockerfile (\texttt{podman build}) y crea el contenedor publicando su puerto. Los scripts de \texttt{defensoria-web}, \texttt{admin-web}, \texttt{revision-web} y \texttt{primer-contacto-web} no incluyen \texttt{-{}-restart always}; se agrega antes de ejecutarlos para que los contenedores vuelvan solos tras reiniciar la VPS:
\begin{lstlisting}
# sed -i 's/podman run -d --name/podman run -d --restart always --name/' \
      $F/*/podman-compose-front*.sh
# grep -h "podman run" $F/*/podman-compose-front*.sh
\end{lstlisting}
Las cinco líneas deben contener \texttt{-{}-restart always}. Después se levantan:
\begin{lstlisting}
# cd $F/front                 && bash podman-compose-front.sh up
# cd $F/front-admin           && bash podman-compose-front-admin.sh up
# cd $F/front-revision        && bash podman-compose-front-revision.sh up
# cd $F/front-primer-contacto && bash podman-compose-front-primer-contacto.sh up
# cd $F/front-subdefensoria   && bash podman-compose-front-subdefensoria.sh up
\end{lstlisting}
\captura{front-03-frontends-up}{Construcción y arranque de los cinco frontends}

\begin{table}[H]
\centering
\caption{Contenedores de los frontends}
\footnotesize
\begin{tabular}{llll}
\toprule
\textbf{Contenedor} & \textbf{Imagen} & \textbf{Puerto VPS $\rightarrow$ contenedor} & \textbf{Ruta pública} \\
\midrule
defensoria-web & defensoria-front-img & 8090 $\rightarrow$ 8090 & \texttt{/} \\
admin-web & defensoria-admin-img & 8091 $\rightarrow$ 80 & \texttt{/admin/} \\
revision-web & defensoria-revision-img & 8092 $\rightarrow$ 80 & \texttt{/revision/} \\
primer-contacto-web & defensoria-primer-contacto-img & 8094 $\rightarrow$ 80 & \texttt{/primer-contacto/} \\
subdefensoria-web & defensoria-subdefensoria-img & 8095 $\rightarrow$ 80 & \texttt{/subdefensoria/} \\
\bottomrule
\end{tabular}
\end{table}

Verificación desde la propia VPS:
\begin{lstlisting}
# podman ps --format "table {{.Names}}\t{{.Status}}\t{{.Ports}}"
# for p in 8090 8091 8092 8094 8095; do
    printf "%s -> " $p; curl -s -o /dev/null -w "%{http_code}\n" http://127.0.0.1:$p/
  done
\end{lstlisting}
Los cinco deben responder \texttt{200}. Desde Internet esos puertos no responden, porque la regla del capítulo \ref{cap:03_seguridad_red_firewall} los cierra; solo se accede a ellos a través de \texttt{router-nginx}.
\captura{front-04-frontends-200}{Contenedores de los frontends en estado Up y respondiendo 200}
